% ---------------------------------------------------------------
%  宏包与记号
% ---------------------------------------------------------------
\usepackage{amsmath,amssymb,amsthm,mathtools}
\usepackage{booktabs}
\usepackage{graphicx}
\usepackage{tikz}
\usetikzlibrary{arrows.meta,calc}
\usepackage{algorithm}
\usepackage{algpseudocode}
\usepackage{listings}
\usepackage{xcolor}
\usepackage{hyperref}
\usepackage{cleveref}

\hypersetup{colorlinks=true,linkcolor=blue!50!black,citecolor=teal!60!black,urlcolor=blue!50!black}

\lstset{basicstyle=\ttfamily\small,keywordstyle=\color{blue!60!black},commentstyle=\color{gray},
  numbers=left,numberstyle=\tiny\color{gray},frame=lines,columns=fullflexible}

% 定理类环境
\theoremstyle{plain}
\newtheorem{theorem}{定理}[section]
\newtheorem{lemma}[theorem]{引理}
\newtheorem{proposition}[theorem]{命题}
\newtheorem{corollary}[theorem]{推论}
\theoremstyle{definition}
\newtheorem{definition}[theorem]{定义}
\newtheorem{example}[theorem]{例}
\theoremstyle{remark}
\newtheorem*{remark}{注}

% 自定义环境: 全文统一的基本假设
\NewDocumentEnvironment{assumption}{o}
  {\par\medskip\noindent\textbf{基本假设\IfValueT{#1}{ (#1)}.}\itshape\ }
  {\par\medskip}

\crefname{theorem}{定理}{定理}
\crefname{lemma}{引理}{引理}
\crefname{proposition}{命题}{命题}
\crefname{corollary}{推论}{推论}
\crefname{definition}{定义}{定义}
\crefname{equation}{式}{式}
\crefname{table}{表}{表}
\crefname{figure}{图}{图}
\crefname{algorithm}{算法}{算法}
\crefformat{section}{第~#2#1#3~节}
\Crefformat{section}{第~#2#1#3~节}
\crefformat{appendix}{附录~#2#1#3}
\Crefformat{appendix}{附录~#2#1#3}

% 记号
\DeclarePairedDelimiter{\norm}{\lVert}{\rVert}
\DeclarePairedDelimiter{\abs}{\lvert}{\rvert}
\DeclarePairedDelimiter{\inner}{\langle}{\rangle}
\newcommand{\R}{\mathbb{R}}
\newcommand{\C}{\mathbb{C}}
\newcommand{\N}{\mathbb{N}}
\newcommand{\BX}{\mathcal{B}(X)}
\newcommand{\KX}{\mathcal{K}(X)}
\DeclareMathOperator{\Ker}{Ker}
\DeclareMathOperator{\Ran}{Ran}
\DeclareMathOperator{\spn}{span}
